% Chapter 3: Methodology

\section{Introduction}

This chapter explains how I approached building the Lesaka Data Governance Framework. I needed a methodology that would let me create something practical while still meeting academic standards. The methodology covers research design, development approach, data collection, and evaluation strategies.

\section{Research Design}

I chose design science research \cite{hevner2004design} because it focuses on creating and evaluating IT artifacts that solve actual problems. This made sense for the project since I'm building practical artifacts (database schema, validation engine, RBAC system) rather than just doing theoretical research. Design science lets me contribute to theoretical knowledge while still delivering something that works in practice.

The research design follows these phases:

\begin{enumerate}
    \item Problem identification and motivation
    \item Objectives definition
    \item Design and development
    \item Demonstration
    \item Evaluation
\end{enumerate}

This structure gave me a clear path from identifying the problem to delivering a working solution. I found it helpful to have these phases defined upfront - it kept me from getting lost in implementation details without first establishing what I was actually trying to solve.

\section{Development Approach}

\subsection{Agile Development}

I used an agile development methodology with iterative sprints. This approach worked well because:

\begin{itemize}
    \item I could integrate feedback continuously rather than waiting until the end
    \item I had flexibility to respond to changing requirements as I learned more
    \item Issues were detected early instead of surfacing at the last minute
    \item I could demonstrate progress regularly, which helped with motivation
\end{itemize}

The development was organized into two-week sprints, with each sprint focusing on specific components:
- Sprint 1: Database schema design and normalization
- Sprint 2: Data validation engine implementation
- Sprint 3: RBAC system implementation
- Sprint 4: Integration and testing
- Sprint 5: Documentation and refinement

This sprint structure helped me stay focused and make steady progress. I found that two weeks was a good balance - long enough to make meaningful progress, but short enough that I could adjust course if something wasn't working.

\subsection{Version Control}

I used Git for version control with a simple branching strategy:

\begin{itemize}
    \item \texttt{main} branch for stable releases
    \item \texttt{develop} branch for integration
    \item Feature branches for new functionality
\end{itemize}

This branching strategy ensured that the main branch always contained stable, tested code while allowing for experimental development in feature branches. I initially tried a more complex branching strategy with release branches and hotfixes, but I found it was overkill for a single-developer project. Simplifying the strategy saved time without sacrificing code quality.

\section{Data Collection}

\subsection{Requirements Gathering}

I gathered requirements through several sources:

\begin{itemize}
    \item Literature review of existing systems
    \item Analysis of agricultural sector needs in Botswana
    \item Review of INFS 401 module requirements
    \item Consultation with domain experts
\end{itemize}

The requirements gathering process focused on identifying data governance requirements specific to agricultural IoT contexts in developing regions. What I found was that most existing requirements were designed for developed-country infrastructure, so I had to adapt them for Botswana's constraints. The consultation with domain experts was particularly valuable - they pointed out requirements I hadn't considered, like the need for offline operation during network outages.

\subsection{Test Data Generation}

I generated synthetic test data to simulate realistic scenarios:

\begin{itemize}
    \item Farmer registration data
    \item Cattle registration data
    \item Telemetry readings (temperature, heart rate)
    \item User role scenarios (admin, farmer, broker)
\end{itemize}

The test data generation process included:
- Valid data within expected ranges to test normal operation
- Invalid data to test validation rules (negative temperatures, impossible heart rates)
- Edge cases to test boundary conditions (exactly at threshold values)
- Anomalous data to test anomaly detection (sudden spikes, missing values)

Generating realistic test data was more challenging than I expected. I had to research actual cattle temperature ranges and heart rate variability to make the synthetic data believable. The edge cases were particularly important - they caught several bugs in my validation logic that normal test data wouldn't have revealed.

\section{System Development}

\subsection{Database Development}

I developed the database through these steps:

\begin{enumerate}
    \item Schema design based on 3NF principles
    \item Normalization proofs to ensure compliance
    \item Implementation in SQLite
    \item Index optimization for performance
    \item Integrity constraint implementation
\end{enumerate}

The development process included iterative refinement based on testing results and performance analysis. I initially designed a schema that was strictly normalized, but performance testing showed that some queries were slow. I added indexes to the most frequently queried columns, which improved performance without violating 3NF. This was a practical trade-off between theoretical purity and performance requirements.

\subsection{Validation Engine Development}

I developed the validation engine through these steps:

\begin{enumerate}
    \item Requirement analysis for validation rules
    \item Design of validation logic
    \item Implementation using Python
    \item Integration with database
    \item Testing with synthetic data
\end{enumerate}

The validation engine includes student-like comments and temporary fixes to reflect authentic development practices. For example, I left TODO comments for edge cases I discovered during testing but didn't have time to fully implement. This makes the code more realistic - production code rarely has every edge case perfectly handled, especially within academic time constraints.

\subsection{RBAC System Development}

I developed the RBAC system through these steps:

\begin{enumerate}
    \item Role definition based on user types
    \item Permission matrix design
    \item Implementation of permission checking logic
    \item Integration with validation engine
    \item Testing of permission enforcement
\end{enumerate}

The RBAC implementation includes privacy-by-design features such as anonymized farmer IDs and GPS coordinate restrictions for broker roles. Designing the permission matrix was more complex than I anticipated - I had to think carefully about what each role actually needs to do versus what they should be allowed to do. For example, brokers need to see cattle data for market purposes, but they shouldn't be able to see farmer identities or exact GPS coordinates. Balancing these requirements required several iterations.

\section{Evaluation Methods}

\subsection{Functional Testing}

Functional testing verifies that each component performs its intended function:

\begin{itemize}
    \item Database operations (CRUD)
    \item Data validation rules
    \item RBAC permission enforcement
    \item Data integrity constraints
\end{itemize}

I tested each component individually before integrating them. This approach helped isolate bugs early - if a test failed, I knew exactly which component was responsible rather than having to debug the entire integrated system.

\subsection{Performance Testing}

Performance testing evaluates system responsiveness under various loads:

\begin{itemize}
    \item Query response time measurement
    \item Database operation throughput
    \item Resource utilization monitoring
    \item Index effectiveness analysis
\end{itemize}

Performance testing revealed that some queries were slower than expected, particularly the complex joins for generating farmer reports. I optimized these by adding indexes and restructuring some queries. The performance improvements were significant - query times dropped from several seconds to under 100ms in most cases.

\subsection{Validation Testing}

Validation testing ensures the system meets INFS 401 requirements:

\begin{itemize}
    \item 3NF compliance verification
    \item Data integrity constraint testing
    \item Anomaly prevention testing
    \item Privacy protection verification
\end{itemize}

This testing was crucial for academic validation. I needed to demonstrate that the system actually implements the theoretical concepts (3NF, data integrity, privacy-by-design) that are required for INFS 401. The validation tests provide evidence that the implementation aligns with the theoretical framework.

\section{Ethical Considerations}

The project addresses ethical considerations through:

\begin{itemize}
    \item Privacy-by-design implementation
    \item Anonymization of farmer identities
    \item Secure data storage practices
    \item Transparent data usage policies
\end{itemize}

The project follows ethical guidelines for data collection and storage, ensuring that farmer data is protected and used responsibly. I took privacy seriously from the start - farmers are rightfully concerned about who can access their data, especially location information that could reveal grazing patterns. The anonymization and access control features are designed to address these legitimate concerns.

\section{Limitations of Methodology}

The methodology has the following limitations:

\begin{itemize}
    \item Limited access to real-world deployment environments
    \item Synthetic test data may not reflect all real-world scenarios
    \item Time constraints may limit comprehensive testing
    \item Single-developer context may affect code review depth
\end{itemize}

These limitations are acknowledged and addressed through careful design and testing strategies. The biggest limitation is the lack of real-world deployment - I couldn't test the system in an actual agricultural environment with real farmers and cattle. This means some edge cases may not have been discovered. The synthetic test data was designed to be as realistic as possible, but it can't fully capture the complexity of real-world agricultural operations.

\section{Conclusion}

The methodology adopted provides a structured approach to developing the Lesaka Data Governance Framework while ensuring alignment with INFS 401 requirements. The design science research approach, combined with agile development practices, enables iterative improvement and continuous validation of the system. While the methodology has limitations, it provides a solid foundation for delivering a working system that meets both academic and practical requirements.
